% 本文件由 LaTeX 排版 Agent 根据有效 translation.json 自动生成。
% author=Augustine of Hippo; work_id=1315; title=论忍耐

\chapter{论忍耐}

% structure: p0001 source=h1 target=omitted reason=page-title-already-rendered-by-parent

\Person{埃拉斯谟}从这篇作品的风格和语言推断，它并非\Person{圣奥古斯丁}所作，并将其与\WorkTitle{论节欲}、\WorkTitle{论爱德之实}、\WorkTitle{论无形之物的信德}等论文归为一类。本笃会编者承认此作具有足以引起怀疑的风格特点（尤其是有意的谐音和押韵结尾，例如第12章：\Quote{\Emph{cautior fuit iste in doloribus quam ille in nemoribus ... consensit ille oblectamentis, non cessit ille tormentis,}}）；然而他们觉得必须依据\Person{奥古斯丁}本人的证言将其保留在其真作之中，\Person{奥古斯丁}在致\Person{达里乌斯}的书信（231, 第7章）[第一卷，第584页]中提到了这篇作品和\WorkTitle{论节欲}。此作未在\WorkTitle{订正录}中被提及，其原因是它似乎曾以讲道的形式发表（见第1章和第3章），而\Person{奥古斯丁}未能实现其进一步撰写一部订正书以回顾其大众讲道和信件的意图（书信224, 第2章）。就内容和教义而言，这篇论文没有任何与\Person{圣奥古斯丁}已知教义和观点相悖或不符之处。\par

1. 那被称作忍耐的心灵之德，是天主如此伟大的恩赐，以至于连赐予我们这恩赐的那一位，其等待恶人悔改的属性，也被称为忍耐（或长久忍耐）。因此，虽然在天主内不可能有痛苦，而\Quote{忍耐}之名源自\Emph{a patiendo}，即来自忍受，但我们不仅忠信地相信一位忍耐的天主，而且有益地宣认祂。然而，天主的忍耐，其性质和程度如何——祂，我们称之为不能受苦的，却并非不耐心，反而极其耐心——谁又能以言语阐述清楚呢？因此，那种忍耐是不可言喻的，正如祂的妒忌、祂的愤怒以及类似的一切。因为如果我们按我们自身的方式来构想这些属性，那么在天主内则根本没有它们。我们体验到这些属性时无不带有烦恼；但绝不可假设天主不能受苦的本性会受到任何烦恼。但正如祂嫉妒而无丝毫心灵昏暗，发怒而无任何扰乱，怜悯而无任何痛苦，后悔而无任何需要纠正的错误；同样，祂忍耐而无任何情感。现在，关于我们能够理解并应当拥有的人类的忍耐，它究竟是怎样的，我将按天主的恩赐和当前论述的简略，尝试阐述。\par

2. 人类的忍耐——是正确、可赞颂并配得上德行之名的——被理解为一种以平静之心容忍恶事的美德，使我们不会以不平静之心弃绝善事，从而能够达到更善的境界。因此，不耐烦的人，由于不愿忍受恶事，非但不能从恶事中得救，反而会遭受更重的恶事。而忍耐的人，宁愿选择不去作恶而忍受，胜过因不忍耐而去作恶，他们既使所忍耐的痛苦变得轻微，也逃脱了因不耐烦而可能陷入的更重恶事。而那些伟大而永恒的善事，他们不会失去，因为他们不屈服于短暂而轻微的恶事：因为\BibleQuote{现今的苦楚与将来在我们身上要显示的光荣，是不能相较的}，如宗徒所说。他又说：\BibleQuote{我们这现时短暂的轻微苦难，正以不可言喻的方式为我们成就永远的伟大光荣}。\par

3. 亲爱的弟兄们，请看人们为了他们所邪恶地爱慕的事物，在劳作和痛苦中忍受何等艰辛，并且他们自认为通过这些事物能变得更加幸福，就更加不幸地贪求它们。为了虚假的财富，为了虚幻的荣誉，为了对游戏和表演的爱好，人们何等耐心地承受极大的危险和麻烦！我们看到那些贪求金钱、荣耀、放纵的人，为了达到他们的欲望，并在得到后不失去它们，他们忍受阳光、雨水、严寒、波浪和最为狂暴的风暴，战争的艰难和不确定性，巨大打击的创伤和可怕的伤口——不是出于无法避免的必然，而是出于应受谴责的意志。但这些疯狂行为在某种程度上被认为是被允许的。因此，贪婪、野心、奢侈以及各种游戏和表演的乐趣，除非他们为了这些而犯下人类法律所禁止的恶行或暴行，否则都被认为属于无辜；不仅如此，一个人若未对任何人造成伤害，无论是为获取或增加金钱，还是为获得或保持荣誉，或是在比赛中竞争、狩猎中，或在掌声中展示某种戏剧场面，而承受了巨大的劳苦和痛苦，不仅不会因世俗的虚荣而受到责备，反而会受到赞扬：因为，正如经上所记载，\BibleQuote{罪人因自己灵魂的欲望而受赞美}。因为欲望的力量使人忍受劳苦和痛苦；除非为了他所享受的事物，没有人会自愿承受令人烦恼的事。但我所说的这些情欲——那些为之燃烧的人为了满足它们而极其耐心地忍受许多艰难和痛苦——被认为是被允许且法律所认可的。\par

4. 不仅如此；难道不正是这样，就连公然的恶行，人们为了施行而非惩罚它们，也忍受了许多最沉重的苦难？世俗文学的作者们不是讲述了一个极其高贵的弑国者，说他能忍受饥饿、干渴、寒冷，这一切他都能够承受，他的身体对缺乏食物、温暖和睡眠的忍耐达到了令人难以置信的地步吗？何必提那些路霸强盗，他们为了埋伏行人，整夜不睡，为了捕捉那些毫无恶意地路过的人，无论天气多么恶劣，都将充满恶念的身心固定在一处？甚至据说他们中有些人习惯于互相轮流折磨，以致这种对抗痛苦的练习和训练与痛苦本身毫无二致。因为，法官折磨他们以通过痛苦之痛获取真相，或许还不如他们自己的同伙折磨他们以通过忍耐痛苦不背叛真相来得厉害。然而，在所有这些情形中，那忍耐与其说值得赞叹，不如说值得谴责：不，既非赞叹也非谴责，因为那并非忍耐；我们只能惊叹其冷酷，否认其忍耐；因为其中没有任何值得正当赞美、值得有益效法的东西；你要正确判断，那灵魂越是将德行的工具交给恶习，它就越是该受更大的惩罚。忍耐是智慧的伴侣，不是私欲偏情的婢女；忍耐是良心的朋友，不是无辜的敌人。\par

5. 因此，当你看见任何人耐心忍受某事时，不要立刻称赞那是忍耐；因为这只能由受苦的原因来证明。当原因良善时，那才是真正的忍耐；当那原因未被私欲玷污时，这忍耐才与虚假有别。但当那原因在于犯罪时，这忍耐在名称上就大错特错了。因为并非所有受苦的人都有份于忍耐，正如并非所有知道的人都有份于知识；那些正确使用受苦的人，才在真实的忍耐中受赞美，才以忍耐的奖赏得冠冕。\par

6. 然而，既然为了私欲甚至恶行，总之为了这暂世的生命和好处，人们竟能奇妙地承受许多可怕苦难的冲击，他们大大地提醒我们，为了善生，该当承受何等伟大的事，好使这生命将来也成为永生，不受时间限制，不损失或浪费任何益处，在真正的幸福中安稳。主说：\Quote{你们要凭着忍耐，保全你们的灵魂。}祂没说：你们的田产、你们的赞美、你们的享乐；而是说：\Quote{你们的灵魂。}既然如此，如果灵魂忍受如此巨大的苦难只为获得那会使它丧亡的事物，那么它该当承受何等大的苦难以免丧亡呢？此外，提一件并非该受指责的事：如果灵魂为了拯救肉体，在医生切割或烧灼的手下忍受如此大的痛苦，那么为了拯救自身，在无论怎样的敌人狂怒之下，它该当承受何等大的痛苦呢？因为医生为使身体不死，用痛苦来为身体的好处着想；而敌人用痛苦和死亡威胁身体，是要催促我们去毁灭灵魂和身体于地狱中。\par

7. 虽然，事实上，如果为了义而轻视身体暂时的生命和好处，并且为了义而最耐心地忍受身体的痛苦或死亡，那么身体的福祉本身反而得到了更明智的眷顾。因为宗徒论及终末的身体的救赎时说：\Quote{我们自己也内心叹息，等待成为义子，即我们身体的救赎。}然后他接着说：\Quote{\InnerQuote{因为我们得救，是出于希望。但看得见的希望不是希望：因为人看见了，为什么还希望呢？}但我们若希望那看不见的，我们就耐心等待。}因此，当任何灾祸确实折磨我们，却不从我们身上榨出恶行时，不仅灵魂因忍耐而得以保全；甚至当身体本身因忍耐而暂时受苦或丧失时，它也被重获，成为永恒稳固和得救，并通过痛苦和死亡为自己积累了不可侵犯的健康和幸福的永生。为此，主耶稣鼓励祂的殉道者忍耐，应许将来身体完全完整，我并非说任何肢体，而是一根头发也不失落。祂说：\Quote{我实在告诉你们，你们头上的一根头发也不会失落。}这样，因为正如宗徒所说，\Quote{从来没有人恨恶自己的肉身}，信友更应以忍耐而非急躁来审慎照料自己肉身的状况，并以将来不朽的无可估量之收益来弥补它现今无论多么重大的损失。\par

8. 然而，忍耐虽是心灵之德，但一部分是心灵在自身内实行之，另一部分是在肉身上实行。在自身内实行忍耐，是指当肉身保持无损无伤时，心灵因事物或言语的任何逆境或污秽而被刺激去行或说某些不合宜或不体面之事，而它耐心地忍受一切恶事，以免自身在行为或言语上犯任何恶。借着这种忍耐，即使我们身体安康，我们仍在这世界的众多冒犯中忍受着我们的福乐被延迟；关于此，有我所引用的那句话说：\BibleQuote{如果我们所希望的尚未看见，我们就借忍耐等待。}\BibleRef{罗 8:25} 借着这种忍耐，圣\Person{达味}忍受了一个辱骂者的斥责；当他本可轻易报复时，他不但没有报复，甚至劝阻了另一个为他恼怒激动的人；他更是以禁止报复而非实施报复来行使他的君王权力。那时他的身体也没有受任何疾病或创伤之苦，而是承认了一个谦卑的时刻，并承行了天主的旨意，为此他以极其忍耐的心饮下了侮辱的苦杯。主也教导了这种忍耐，当仆人们因稗子混入而激动，想要收集它们时，祂答说：\BibleQuote{让两者一起生长直到收获。}\BibleRef{玛 13:30} 也就是说，必须忍耐那不应急于除掉的。主亲自作出了这种忍耐的榜样，在祂肉身受难之前，祂如此忍耐了祂的门徒\Person{犹达斯}，以致在指出他是叛徒之前，祂容忍他作贼；在经历捆绑、十字架和死亡之前，祂对那些充满诡诈的嘴唇并没有拒绝平安之吻。所有这些以及其它懒于例举的事，都属于那种忍耐：心灵并非自身之罪，而是来自外部的任何恶事，在肉身完全无损的情况下，耐心地在自身内忍受。但另一种忍耐是同一心灵在肉身受苦时承受任何烦恼和痛苦；不像愚昧或邪恶之人为了获取虚浮之物或行凶作恶而如此，而是如主所说，\BibleQuote{为了义德的缘故}\BibleRef{玛 5:10}。在两种忍耐中，圣殉道者都曾奋力争战。因为他们既充满了不敬神之人的轻蔑责难——此时肉身未受损伤，心灵却有其（可以说是）打击和创伤，并且他们坚忍不拔地忍受这些；同时在身体上他们被捆绑、监禁、受饥渴之苦、被折磨、被割伤、被撕裂、被焚烧、被屠杀；他们以不可动摇的虔诚将心灵交托于天主，而肉身则忍受着刁钻残忍所能设想的一切痛苦。\par

9. 这确实是一场更大的忍耐之战，当不再是可见的敌人借着迫害和狂怒诱使我们犯罪——这样的敌人可以在光天化日之下以不予同意而战胜；而是魔鬼本身（他也借着不信之徒之子，如同他的器皿，迫害光明之子）亲自隐蔽地攻击我们，以他的狂怒怂恿我们去做或说一些反对天主的事。正如圣\Person{约伯}所经历的，他两种诱惑都受折磨，但靠着忍耐的坚定力量和虔诚的武器，在两者中都未被打败。因为首先，在身体未受损伤的情况下，他失去了一切所有，为的是使心灵在肉身受折磨之前，因失去人们通常极为珍视的事物而崩溃，并因失去那些他本以为敬拜天主所依靠的事物而说出反对天主的话。他也突然遭受全部儿子丧亡的打击，以致他一个接一个生育的儿女，一下子全部失去，仿佛他们的众多不是为他幸福增添光彩，而是为了增加他的灾祸。但当他承受了这些，仍坚定不移地依靠他的天主时，他紧附于祂的旨意，而这旨意除非他自愿，否则不可能失去；并且他找到那永远不会失去的祂，代替他所失去的一切。因为拿走这些的并非那个有意伤害的仇敌，而是那赐予仇敌伤害能力的祂。接着，仇敌也攻击身体，现在不是攻击人外在的一切，而是攻击人自身，凡他能触及的部分，他都击打。从头到脚是灼热的痛苦，是蠕动的蛆虫，是流脓的疮口；然而腐烂的身体中心灵仍保持完整，并且尽管这腐烂肉身带来恐怖的折磨，它以不可侵犯的虔诚和无玷的忍耐忍受了一切。他的妻子站在一旁，非但没有给她丈夫任何帮助，反而怂恿他亵渎天主。因为我们不应认为魔鬼在使她留存活下来而夺走儿子时，是一个不善作恶的工匠；相反，他在\Person{厄娃}身上早已学会了她对诱惑者是多么必需。但现在他找不到第二个\Person{亚当}，能借着女人俘获他。\Person{约伯}在他悲伤的时刻比\Person{亚当}在他喜悦的凉亭中更为谨慎；一个在愉快事物中被打败，另一个在痛苦中得胜；一个向似乎可喜之事屈服，另一个在最可怕的折磨中毫不退缩。他的朋友们也站在一旁，不是在他苦难中安慰他，而是怀疑他作恶。因为当他承受如此巨大的痛苦时，他们不相信他是无辜的，他们的舌头也不禁说出他良心没有说过的话；这样，在肉身无情的折磨中，他的心灵也遭到不实指责的打击。但他以肉身承受自身的痛苦，以内心承受他人的错误，斥责妻子的愚妄，教导友人智慧，在每一件事和每一个人身上保持了忍耐。\par

10. 这些人正该受那些因欲得活命而受人搜捕时便自寻短见的人的注目；他们既剥夺了眼前的生命，也拒绝了那将要来的。若有人迫使他们否认基督或行任何违背正义之事，他们本当如真殉道者那样，耐心忍受一切，而不应鲁莽地自寻死路。倘若以逃避苦难为由而自杀是合理的，那么圣约伯大可自杀，好从魔鬼的残酷所加于他的财产、儿子和肢体上的巨大苦难中逃脱。但他并没有这样做。一个智慧人，决不会对自己做出连那愚妇也未曾建议的事。而倘若她曾建议过，那她在此也必当得到她建议亵渎时所受的答复：\BibleQuote{你说话像愚顽的妇人一样。我们从主的手里得福，难道不也受害吗？}因为连他也将失去忍耐，倘若或照她所建议的亵渎，或照她连提都不敢提的自杀，他死了，便列在那些所记的\BibleQuote{祸哉，那失去忍耐的人！}之中，并且痛苦不但没有减轻，反而增加，因为在肉身死后，他将被拖去受惩罚，无论是亵渎者的、杀人者的、还是比杀亲者更恶的人的惩罚。既然杀亲者之所以比任何杀人者更邪恶，是因为他杀的不单是一个人，而且是一个亲属；而在杀亲者中，被杀者越亲近，其罪越大；那么，毫无疑问，自杀者更为恶劣，因为没有人比他自己更亲近自己了。这些可怜的人究竟是什么意思？他们既在此世加痛苦于自身，又因对天主的亵渎及对自己的残忍来世理当承受祂所加的痛苦，竟还寻求殉道者的光荣？即使他们因基督的真见证而受迫害，并为了免遭迫害者的任何痛苦而自杀，也可正当地对他们说：\BibleQuote{祸哉，那失去忍耐的人！}因为若连不耐心的痛苦也获得冠冕，忍耐又怎能获得公义的赏报？或者，那个被称为\BibleQuote{你应爱你的近人如你自己}的人，若他对自己行了禁止对近人行的谋杀，又怎能被判定为无罪呢？\par

11. 愿圣徒们从圣经中聆听忍耐的训诫：\BibleQuote{我儿，你前去事奉天主时，应站在正义与敬畏中，预备你的灵魂受试探；谦卑你的心，坚忍到底；使你的生命在末后得以增长。凡临到你的事，你都要接受；在痛苦中要坚忍，在卑微中要有忍耐。因为金银在火中试炼，而蒙悦纳的人在谦卑的炉中试炼。}在另一处我们读到：\BibleQuote{我儿，不可轻视上主的惩戒，也不可因祂的责备而厌倦；因为上主所爱的，祂必惩戒，祂鞭打所收纳的每一个儿子。}这里所说的\Quote{所收纳的儿子}，与上述证言中的\Quote{蒙悦纳的人}是相同的。因为这是公义的：我们因悖逆地贪求享乐之物而被逐出起初的乐园福乐，却要通过谦卑忍耐烦恼之事而被接纳回来；因作恶而被驱逐，因受苦而回归；在那里违背正义而行恶，在这里为正义而忍耐苦难。\par

12. 但关于真正的忍耐，即配得上这一美德之名的忍耐，它从何而来，现须加以探究。因为有些人将其归诸人的意志之力——这力并非天主助佑所赐，而是来自自由意志。这种错误乃是骄傲的错误：因为这是那些富有之人的错误，经上论及他们说：\BibleQuote{富足的人受讥讽的责斥，骄傲的人受轻视。}因此，并非那种\BibleQuote{永不消灭的穷人的忍耐。}因为这些穷人从那位富足者——对祂曾说过\Quote{祢是我的天主，因祢不需要我的财物}，一切\BibleQuote{美好的恩赐和完美的赠予}都从祂而来——那里领受忍耐；困苦贫穷的人向祂呼求，在祈求、寻找、敲门时说：\BibleQuote{我的天主，求祢救我脱离罪人的手，脱离无法无天和不义之人的手；因为祢是我的忍耐，上主，我自幼以来的盼望。}但这些富有的人，自满自足，不屑在天主面前成为有缺乏的，以免从祂那里领受真正的忍耐；他们夸耀自己虚假的忍耐，试图\Quote{扰乱穷人的计谋，因为上主是他的盼望。}他们也不顾自己只是凡人，竟如此依赖自己的，即人的意志，以致陷入所记的\BibleQuote{凡信赖世人的人，是可诅咒的。}因此，即便他们偶然在艰难困苦中坚忍下来，无论是为了不使人不悦，还是为了避免更糟的事，或是出于自满和自恃的骄傲意志，那么关于忍耐，便可以对他们说那蒙福的宗徒雅各伯关于智慧所说的话：\BibleQuote{这种智慧不是从上而来的，而是属地的、属血肉的、属魔鬼的。}既然骄傲的人可以有虚假的智慧，为何不能有虚假的忍耐呢？真正的智慧从何而来，真正的忍耐也从何而来。因为那神贫的人向祂歌唱：\BibleQuote{我的灵魂臣服于天主，因为我的忍耐来自于祂。}\par

13. 但他们回答说：\Quote{如果人的意志，在没有天主任何帮助的情况下，凭自由抉择的力量，为了享受这必死生命和罪恶的欢愉，而承受如此众多严重可怕的痛苦（无论是精神还是肉体上的），那么，为何同一个人意志不能以同样的方式，凭同样的自由抉择的力量，不求天主的帮助，而是依靠自身自然能力的充足，为了正义和永生，在它所承受的一切劳苦或忧伤中，最坚忍地承受同样的事呢？或者，他们说，不义者的意志，没有天主的帮助，就足以使他们甚至自我折磨，为了不义而经受酷刑，并在被他人折磨之前就足以如此？那些贪爱今生喘息机会者的意志，没有天主的帮助，就足以在最残酷持久的痛苦中坚持谎言，以免承认罪行而被处死？而义人的意志，除非从上获得力量，否则就不足以忍受任何痛苦，无论是为了正义本身之美，还是为了永生之爱？}\par

14. 说这些话的人不明白：每个恶人之所以更能忍受任何苦难，是因为他心中的世俗贪欲更强大；同样，每个义人之所以更能勇敢忍受任何苦难，是因为他心中天主的爱更强大。但世俗的贪欲始于意志的选择，进于快乐的满足，固于习惯的锁链；而\BibleQuote{天主的爱已倾注在我们心中}，这并非来自我们自己，而是\BibleQuote{藉着所赐予我们的圣神}。因此，义人的忍耐来自那一位，藉着祂，他们的爱（对祂的爱）被倾注。宗徒在赞扬和显扬这爱（爱德）时，在它其它的美善品质中，说它\BibleQuote{含忍一切}。他说：\Quote{\InnerQuote{爱德是宽宏的}}。稍后他又说：\BibleQuote{忍受一切}。所以，圣人们心中天主的爱（或仁爱）越大，他们就越为他们所爱的那位忍受一切；而罪人们心中的世俗贪欲越大，他们就越为他们所贪恋的对象忍受一切。因此，义人的真正忍耐来自同一天主的爱，而不义之人的虚假忍耐来自同世俗的贪欲。关于这一点，若望宗徒说：\BibleQuote{不要爱世界，也不要爱世界上的事；谁若爱世界，父的爱就不在他内；因为凡世界上的事，就是肉身的贪欲、眼目的贪欲、以及人生的骄矜，都不是出自父，而是出自世界。}因此，这种不是出自父而是出自世界的私欲偏情，在一个人的心中越激烈炽热，他就越能忍受一切痛苦和忧患，为的是他所贪恋的。所以，如上所述，这不是从上而来的忍耐，而虔敬之人的忍耐是从上而来的，从光明之父降下的。因此，那个是属地的，这个是属天的；那个是属血肉的，这个是属神的；那个是属魔鬼的，这个是属天主的。因为私欲偏情——罪人们正是由于它而顽固地忍受一切——是出自世界的；而爱德——正人君子正是由于它而勇敢地忍受一切——是出自天主的。因此，对于那种虚假的忍耐，人的意志可能无需天主的帮助就足够了；它越是渴望情欲，就越是坚硬，并且因自身变得越恶劣而更持久地忍受苦难：而对于这种真正的忍耐，人的意志除非从上得到帮助和燃烧，否则就不足以承受苦难，因为圣神是它的火；除非由祂点燃对那不可变的善的爱，否则就不能忍受所承受的恶。\par

15. 因为，正如神圣的证言所证明：\BibleQuote{天主是爱，那存留在爱内的，就存留在天主内，天主也存留在他内。}所以，凡主张没有天主的帮助也能拥有天主的爱，他除了主张没有天主也能拥有天主之外，还能主张什么呢？哪个基督徒会这样说呢？连疯子也不敢这样说。因此，在宗徒中，那真实的、虔敬的、忠信的忍耐，藉着圣人们的口，喜乐地说：\BibleQuote{谁能使我们与基督的爱隔绝呢？是困苦吗？是窘迫吗？是迫害吗？是饥饿吗？是赤贫吗？是危险吗？是刀剑吗？正如经上所载：\InnerQuote{为了你，我们整天被置于死地，被视作待宰的羊。}然而，靠着那爱我们的主，我们在这一切事上大获全胜：}不是靠我们自己，而是\BibleQuote{靠着那爱我们的主}。然后他接着说：\BibleQuote{因为我深信：无论是死亡，是生命，是天使，是掌权者，是现存的，是未来的，是有能的，是高处，是深处，或是任何其它受造物，都不能使我们与天主的爱相隔绝，这爱是在我们的主基督耶稣内的。}这就是那\BibleQuote{天主的爱}，是\BibleQuote{藉着所赐予我们的圣神，倾注在我们心中的}。但恶人的私欲偏情，使他们有虚假的忍耐，\BibleQuote{不是出自父}，正如若望宗徒所说，而是出自世界。\par

16. 这里有人会说：\\Quote{如果恶人的私欲偏情（由此他们为了所贪恋的事物忍受一切邪恶）是属于世界的，那么它怎么被称为出于他们自己的意志呢？} 仿佛他们自己在爱世界、离弃那创造世界的主时，不也是属于世界似的。因为 \\BibleQuote{他们侍奉受造物甚于造物主，祂是永受赞美的。} 所以，无论宗徒\\Person{若望}用\\Quote{世界}一词是指爱世界的人，那出于他们自己的意志便是属于世界的；还是他用世界的名称涵盖天地及其中所有的一切，即普遍受造物，显然这受造物的意志，既然不是造物主的，就是属于世界的。因此，主对这样的人说：\\BibleQuote{你们是出于下，我是出于上：你们属于这个世界，我不属于这个世界。} 祂又对宗徒说：\\BibleQuote{如果你们属于世界，世界会爱自己的。} 但唯恐他们过于自恃，当祂说他们不属于世界时，他们以为这是出于本性，而非出于恩宠，因此祂说：\\BibleQuote{但因为你们不属于世界，而是我从世界中拣选了你们，所以世界恨你们。} 由此可见，他们曾经是属于世界的；因为为使你们不属于世界，你们是从世界中被拣选出来的。\par

17. 现在宗徒\\Person{保禄}证明这拣选不是由于先行的善行功绩，而是恩宠的拣选，如此说：\\BibleQuote{在现今时期，因恩宠的拣选，留下了一部分余民。但假使是出于恩宠，便不是出于行为；否则恩宠就不再是恩宠了。} 这就是恩宠的拣选；即藉着天主的恩宠，人得以被选。我说，这是先于一切人的善行功绩的恩宠拣选。因为如果它是给予任何善行功绩的，那么它就不再是无偿的给予，而是作为债务偿还，因而并非真正的恩宠；如同同一宗徒\\Person{保禄}说：\\BibleQuote{报酬不算是恩宠，而是债务。} 然而，如果为了使它是真正的恩宠，即无偿的，它在人内找不到任何应得功绩的事（这在以下经文中得到很好理解：\\BibleQuote{祢要白白地拯救他们}），那么它本身必赐予功绩，而不是给予功绩。因此，它甚至先于信德，而一切善行都是从信德开始的。\\BibleQuote{因为义人}，正如经上所载，\\BibleQuote{因信德而生活。} 而且，恩宠不仅帮助义人，也成就不虔敬的人。因此，即使当它帮助义人，似乎是对其功绩的回报，即便如此它也不停止是恩宠，因为它所帮助的，正是它自己所赐予的。因此，鉴于这先于一切人的善行功绩的恩宠，基督不仅被不虔敬的人处死，而且\\BibleQuote{为不虔敬的人死了。} 在祂死之前，祂拣选了宗徒们，当然不是那时已是义人，而是将要成为义人：祂对他们说：\\BibleQuote{我已从世界中拣选了你们。} 因为祂对谁说：\\BibleQuote{你们不属于世界}，然后唯恐他们认为自己从未属于世界，随即补充道：\\BibleQuote{但我从世界中拣选了你们}；显然，他们不属于世界是出于祂自己的拣选所赐予他们的。因此，如果他们是因自己的义德而不是因祂的恩宠被拣选，他们就不会从世界中被选出来，因为如果他们已经是义人，他们就已经不属于世界了。再者，如果他们被拣选的原因是他们已经是义人，那么他们就已经先拣选了主。因为谁能是义人而不选择义德呢？\\BibleQuote{但是律法的终极是基督，使一切相信的人都获得义德。祂为我们成了天主智慧、义德、圣化和救赎：正如经上所载：夸耀的，当在主内夸耀。} 祂自己就是我们的义德。\par

18. 因此，古时的义人，在圣言降生成人之前，藉着这基督的信德和这真实的义德（基督为我们成了义德），得以成义；他们相信那将要来临的，正如我们相信那已经来临的。他们自己也是藉着恩宠，因信德而得救，不是出于自己，而是天主的恩赐，不是出于行为，免得有人自夸。因为他们的善行并不先于天主的怜悯，而是跟随怜悯。因为远在基督降生成人之前，他们就听到并写下：\\BibleQuote{我要怜悯谁就怜悯谁，我要对谁动慈心就对谁动慈心。} 从这些天主的话语，宗徒\\Person{保禄}在很久以后说：\\BibleQuote{如此看来，这不在于那意愿的，也不在于那奔跑的，而在于那显示怜悯的天主。} 这也是他们自己的声音，远在基督降生成人之前：\\BibleQuote{我的天主，祢的怜悯必要先迎我。} 他们怎么可能与基督的信德无关呢？正是藉着他们的爱德，基督甚至预先被宣告给我们；没有祂的信德，任何一个有死的人，无论过去、现在或将来，都不能成为义人。如果宗徒们当时已经是义人而被基督拣选，那么他们就会先拣选祂，以便义人能被他拣选，因为没有祂他们不能成为义人。但事实并非如此：正如祂亲自对他们说：\\BibleQuote{不是你们拣选了我，而是我拣选了你们。} 关于这一点，宗徒\\Person{若望}说：\\BibleQuote{不是我们爱了天主，而是祂爱了我们。}\par

19. 既然情况如此，人在此世使用自己的私意时，在他选择并爱慕天主之前，难道不都是不义和不敬的吗？我说：\Quote{人是什么？}一个背离造物主的受造物，除非他的造物主\Quote{眷顾他}，并自由地选择他、自由地爱慕他？因为他自己无法选择或爱慕，除非他先被选择、被爱慕并得到治愈，因为他在选择盲目时便看不见，在爱慕懈怠时便很快厌倦。但或许有人会说：天主如何首先选择并爱慕不义的人，以使他们成义呢？经上记载：\BibleQuote{上主，祢恼恨所有作恶的人。}我们认为，这是以何等奇妙而不可言喻的方式？然而，我们也能想象，良医既恨又爱病人：因他患病而恨他，因要驱除他的疾病而爱他。\par

20. 关于爱德，我们已说了这些。若没有爱德，我们里面便不可能有真正的忍耐，因为在善人身上，是天主的爱忍受一切，而在恶人身上，则是世界的贪欲。但这爱在我们内，是藉着赐予我们的圣神而来的。因此，爱由谁而来，忍耐也由谁而来。然而，世界的贪欲，当它耐心忍受任何灾祸的重担时，便夸耀自己意志的力量，如同夸耀疾病的麻木，而非健康的强韧。这种夸耀是狂妄的：这不是忍耐的言语，而是昏聩的言语。这样的意志，越是贪婪现世的善，便越能忍受痛苦的恶，因为它更缺乏永恒的事物。\par

21. 但如果这意志被魔鬼的灵以欺骗的景象和不洁的刺激所驱使和点燃，与其狼狈为奸地联合共谋，这灵便使人意志要么因错误而疯狂，要么因对某种世俗快乐的贪欲而燃烧；因此，它似乎显示出对难以忍受的苦难有惊人的耐力：然而，这并不意味著没有其他不洁之灵的煽动，一个恶意的意志就不能存在，如同一个善意的意志没有圣神的帮助就不能存在一样。因为即使没有任何诱惑或煽动的灵，恶意的意志也能存在，这在魔鬼本身身上就足够清楚了：他之所以成为魔鬼，并非通过其他魔鬼，而是出于他自己的意志。因此，一个恶意的意志，无论是被贪欲催促，还是被恐惧召回，或是被喜乐扩张，或是被悲伤收缩，并在所有这些心灵的扰动中忍耐并轻视那些对他人或另一个时刻更为痛苦的事，这恶意的意志无需其他灵来驱使，便能自我诱惑，并从高处堕落至低处。它把那东西——无论是它寻求获得或害怕失去的，或是因得到而欢喜或因失去而悲伤的——看得越愉悦，便越能为了它而忍受稍微不那么痛苦的事。因为无论那东西是什么，都是受造物，其乐趣是可知的。因为在某种意义上，所爱的受造物接近那爱它的受造物，以亲密的接触和连接，向其体验自己的甜蜜。\par

22. 但造物主的乐趣却远非如此，正如经上所记：\BibleQuote{也要祢快乐之河的赐予他们喝。}它不是像我们一样的受造物。因此，除非从那里赐予我们祂的爱，否则我们里面便没有它的源头。因此，一个善意的意志，即我们藉以爱天主的意志，除非天主也在其中工作，使人愿意，否则不能在人身内。所以，这善意的意志——即忠信服从天主的意志，被来自上界的圣洁热情点燃的意志，爱天主并为天主而爱邻人的意志——无论是出于爱，如宗徒伯多禄所说：\BibleQuote{主，祢知道我爱祢；}或是出于恐惧，如宗徒保禄所说：\BibleQuote{怀着恐惧和战栗，成就你们的得救；}或是出于喜乐，如他所说：\BibleQuote{在希望中喜乐，在磨难中忍耐；}或是出于悲伤，如他说他因弟兄们心中多有忧愁；无论它以何种方式忍受怎样的痛苦和艰难，都是天主的爱\BibleQuote{忍受一切}，这爱藉着赐予我们的圣神倾注在我们心中。为此，虔敬毫不怀疑，正如圣洁之人的爱德一样，虔诚忍耐之人的忍耐也是天主的恩赐。因为神圣的圣经既不能欺骗也不被欺骗，它不仅旧约中有此事的见证，如向天主所说：\BibleQuote{祢是我的忍耐}，以及\BibleQuote{我的忍耐由祂而来；}又如另一位先知所说，我们领受了刚毅之神；而且在宗徒著作中我们也读到：\BibleQuote{因为为了基督，赐给你们的不但是信祂，而且也为祂受苦。}因此，人被告知他因别人的慈悲而受惠时，不应让那使心意因自己的功劳而高傲。\par

23. 此外，如果有人没有爱德——这爱德属于圣神的合一与和平的连系，天主教会藉此聚集并联络在一起——却陷入某种分裂，为了不否认基督而忍受患难、困苦、饥饿、赤身露体、迫害、危险、监禁、捆锁、酷刑、刀剑、火焰、野兽，甚至十字架，出于对地狱和永火的恐惧：这决不应受责备，反而这也是一种值得称赞的忍耐。因为我们不能说，他若否认基督而不受这些苦，就比他承认基督而受苦更好；相反，我们必须认为，在审判时，这或许会比否认基督而避免这一切更可容忍。正如宗徒所说：\BibleQuote{我若把我所有的财物施舍给人，又舍身被火焚烧，却没有爱德，为我毫无益处。}应理解为对获得天国毫无益处，但并非对于在最后审判时承受较轻的惩罚毫无益处。\par

24. 然而，我们很可能要问：这种忍耐是否也是天主的恩赐，还是应归因于人类意志的力量？藉此忍耐，一个与教会分离的人，并非为了使他分离的错误，而是为了仍与他同在的圣事或圣言的真理，出于对永恒痛苦的恐惧而忍受暂时的痛苦。因为我们必须小心，免得我们若声称那忍耐是天主的恩赐，拥有它的人就被认为也属于天主的国；但若我们否认它是天主的恩赐，我们就不得不承认，在没有天主的助佑和恩赐的情况下，人的意志中也能有某种善。因为不可否认，一个人若相信自己若否认基督将遭受永恒惩罚的痛苦，并为此信德忍受并轻视人所施加的任何惩罚，这是一件善事。\par

25. 同样，既然我们不否认这是天主的恩赐，我们必须这样理解：有一些天主的恩赐是属于那在上耶路撒冷、自由、为众人之母的儿女的（因为这些在某种程度上是遗产，在其中我们\BibleQuote{是天主的继承者，与基督同为继承者}：但另有一些恩赐，甚至连妾的儿女也可以领受，那些肉身的犹太人和分裂者或异端者被比作他们。因为经上记载：\BibleQuote{把婢女和她的儿子赶出去，因为婢女的儿子不能与我的儿子依撒格一同继承产业}：天主曾对亚巴郎说：\BibleQuote{在依撒格内，你的后裔将被称}：宗徒将此解释为：\BibleQuote{这就是说，不是肉身所生的儿女是天主的儿女，而是应许的儿女才被认为是后裔}，这使我们明白，亚巴郎的后裔因基督而属于天主的儿女，他们是基督的身体和肢体，即唯一、真实、独生、至公的天主教会，持守着虔诚的信德；不是那藉着高傲或恐惧而运作的信德，而是\BibleQuote{藉着爱德运作的}信德；然而，当亚巴郎将妾的儿女从他儿子依撒格那里打发走时，他也没有忽略赐给他们一些礼物，使他们不至于完全空手离去，但并非使他们成为继承者。因为如此我们读到：\BibleQuote{亚巴郎将一切所有都给了依撒格；又将礼物赐给妾的众子，打发他们离开他的儿子依撒格}。那么，如果我们自由耶路撒冷的儿女，就让我们明白，被排除在遗产之外者的礼物与继承者的礼物是不同的。因为继承者就是那些被告知\BibleQuote{你们所领受的不是奴役的灵，以致再度恐惧，而是领受了作子女的灵，藉此我们呼叫：阿爸，父啊}的人。\par

26. 因此，让我们以爱德的精神呼喊，直到我们进入那要永远安息的遗产；让我们藉着那适合自由人的爱德，而非适合奴隶的恐惧，忍耐痛苦。让我们呼喊，只要我们仍是贫穷的，直到我们与那遗产一同致富。我们已领受了如此大的抵押，因基督为使我们致富而成为贫穷；祂被高举至天上的财富时，那应吹入我们心中神圣渴望的圣神被派遣来了。这些贫穷者，目前是相信而非看见，是希望而非享受，是叹息渴望而非在福乐中为王，是饥渴而非满足的；那么，这些贫穷者的\BibleQuote{忍耐不会永远消失}：并非在那里也有忍耐，因为那里将没有可忍耐的事；而是\BibleQuote{不会消失}，意思是它将不会没有果实。它的果实将永远持续，所以\BibleQuote{不会永远消失}。因为那徒劳劳动的人，当他的希望落空时，有理由说：\Quote{我浪费了这么多劳力}；但那达到其劳动应许的人，会自我庆贺说：\Quote{我没有浪费我的劳动}。因此，劳动被称为不消失（或损失），并非因为它持续永远，而是因为它没有白费。同样，基督的贫穷者（他们作为基督的继承者将致富）的忍耐不会永远消失：并非在那里我们也将被命令忍耐，而是因为我们在此所忍耐的，将使我们享受永恒的福乐。祂将不终止永恒的福乐，祂给予意志以暂时的忍耐：因为两者都是祂赐予爱德的恩赐，而爱德本身也是祂的恩赐。\par

\SourceRecord{Translated by H. Browne.；Nicene and Post-Nicene Fathers, First Series；Vol. 3.；Edited by Philip Schaff.；Buffalo, NY: Christian Literature Publishing Co.,；1887.；<http://www.newadvent.org/fathers/1315.htm>.}
